% Source presentations for register row N-c-021.

\subsection{The relational, closure, completeness, and complexity presentations of PDL}
\label{sec:p09-presentation-n-c-021}

\subsubsection{Pratt's bridge from triples to internal modalities}

Let \(S\) be a nonempty set of states, let a state assertion be a subset of
\(S\), and let the program \(a\) denote a relation \(R_a\subseteq S\times S\).
Pratt's external partial-correctness notation has the exact satisfaction
condition
\[
 P\{a\}Q
 \quad\Longleftrightarrow\quad
 \forall s,t\in S\,
   \bigl(s\in P\land (s,t)\in R_a\;\Longrightarrow\;t\in Q\bigr).
\]
The internal modal formulas instead have a truth value at one state:
\[
\begin{aligned}
 s\models[a]Q
   &\quad\Longleftrightarrow\quad
   \forall t\in S\bigl((s,t)\in R_a\Longrightarrow t\in Q\bigr),\\
 s\models\langle a\rangle Q
   &\quad\Longleftrightarrow\quad
   \exists t\in S\bigl((s,t)\in R_a\land t\in Q\bigr).
\end{aligned}
\]
Thus programs, external ordered pairs of assertions, and internal modal
formulas have three different standings.  Expanding the two definitions proves
the bridge, rather than merely suggesting it:
\[
 P\{a\}Q
 \quad\Longleftrightarrow\quad
 \forall s\in S\,(s\in P\Longrightarrow s\models[a]Q)
 \quad\Longleftrightarrow\quad
 P\to[a]Q\text{ is valid}.
\]

The regular operations are already visible at the relational level.  With
composition read from left to right,
\[
\begin{aligned}
 R_{a;b}
   &=\{(s,t):\exists u\,((s,u)\in R_a\land(u,t)\in R_b)\},\\
 R_{a\cup b}&=R_a\cup R_b.
\end{aligned}
\]
Direct substitution in the modal clauses gives
\[
\begin{aligned}
 s\models[a;b]Q
   &\Longleftrightarrow s\models[a][b]Q,\\
 s\models[a\cup b]Q
   &\Longleftrightarrow s\models[a]Q\land[b]Q,\\
 s\models\langle a\cup b\rangle Q
   &\Longleftrightarrow
     s\models\langle a\rangle Q\lor\langle b\rangle Q,\\
 s\models[a]Q
   &\Longleftrightarrow s\models\neg\langle a\rangle\neg Q.
\end{aligned}
\]
For example, take \(S=\{s,t\}\), \(P=\{s\}\), \(Q=\{t\}\),
\(R_{a_{\mathrm{halt}}}=\{(s,t)\}\), and
\(R_{a_{\mathrm{div}}}=\varnothing\).  Both
\(P\{a_{\mathrm{halt}}\}Q\) and \(P\{a_{\mathrm{div}}\}Q\) hold: the second
does so vacuously because it has no final state.  At \(s\), however,
\[
 \langle a_{\mathrm{halt}}\rangle\top
 \qquad\text{and}\qquad
 \neg\langle a_{\mathrm{div}}\rangle\top
\]
are true.  This pair isolates exactly what the external partial-correctness
assertion cannot express: existence of a terminating transition
\cite[technical-report pp.~4--10]{Pratt1976}.

Pratt proves semantic results for the operations and fragments he actually
defines.  The report predates the later Fischer--Ladner closure, the
Kozen--Parikh completeness proof, and the final EXPTIME classification; those
results are not back-attributed to it.

\subsubsection{Fischer--Ladner closure and the finite-model argument}

Definition~\ref{def:p09-pdl} gives the mutually generated formula and program
syntax: the mutuality is required because a formula may contain a
program-indexed modality and a program may contain a formula test.  In a model
\(M=(S,(R_a)_a,V)\), the compound relations are
\[
\begin{aligned}
 \llbracket a\rrbracket_M&=R_a,\\
 \llbracket\alpha\cup\beta\rrbracket_M
   &=\llbracket\alpha\rrbracket_M\cup\llbracket\beta\rrbracket_M,\\
 \llbracket\alpha;\beta\rrbracket_M
   &=\{(s,t):\exists u\,
       ((s,u)\in\llbracket\alpha\rrbracket_M\land
        (u,t)\in\llbracket\beta\rrbracket_M)\},\\
 \llbracket\alpha^{*}\rrbracket_M
   &=\bigcup_{n<\omega}\llbracket\alpha\rrbracket_M^{\,n},\\
 \llbracket B?\rrbracket_M
   &=\{(s,s):M,s\models B\}.
\end{aligned}
\]
Here the zeroth relational power is the identity relation.  Formula truth is
Boolean truth together with
\[
 M,s\models[\alpha]B
 \quad\Longleftrightarrow\quad
 \forall t\bigl((s,t)\in\llbracket\alpha\rrbracket_M
                 \Longrightarrow M,t\models B\bigr),
\]
and \(\langle\alpha\rangle B\) is the dual
\(\neg[\alpha]\neg B\).  These clauses separately define the standing of a
program term, a formula, and the metalinguistic relation denoted by a program
\cite[\S\S1--2]{FischerLadner1979}.

Ordinary subformulas are not enough for filtration, because a box indexed by a
compound program owes the one-step formulas through which its obligation is
propagated.  In the Boolean basis fixed by this handbook, the exact simultaneous
recursion for \(\operatorname{FL}\) and
\(\operatorname{FL}_{\Box}\) is printed immediately before this presentation.
It is the modern box-form normalization of the Fischer--Ladner construction;
the auxiliary function decomposes only the outer program and thereby makes the
star clause well-founded \cite[\S6.1]{HarelKozenTiuryn2000}.

For a complete calculation, put
\[
 C_0=\left[(a;b^{*})\cup(p?;c)\right]q.
\]
Repeated application of the union, composition, star, test, and atomic-program
clauses gives exactly
\[
\begin{aligned}
 \operatorname{FL}(C_0)=\{\,
   &C_0,
   [a;b^{*}]q,
   [a][b^{*}]q,
   [b^{*}]q,
   [b][b^{*}]q,\\
   &[p?;c]q,
   [p?][c]q,
   p,
   [c]q,
   q
 \,\}.
\end{aligned}
\]
The entries \([b][b^{*}]q\) and \(p\) are not ordinary subformulas of
\(C_0\): the first is owed by star unfolding and the second by evaluation of the
test.  Conversely, this ten-formula set does not close under arbitrary program
or formula formation.

Here is the finite-model proof at the level needed to see why those extra
entries suffice.  Let \(C\) be satisfiable in \(M\), put
\(\Delta=\operatorname{FL}(C)\), and define
\[
 s\equiv_{\Delta}t
 \quad\Longleftrightarrow\quad
 \forall B\in\Delta\,
   (M,s\models B\Longleftrightarrow M,t\models B).
\]
There are at most \(2^{|\Delta|}\) equivalence classes.  On those classes define
the filtered atomic edge by
\[
 [s]R_a^{f}[t]
 \quad\Longleftrightarrow\quad
 \exists s'\equiv_{\Delta}s\,
 \exists t'\equiv_{\Delta}t\,((s',t')\in R_a),
\]
and interpret the compound programs from these atomic edges by the five
relational clauses above.  Propositional letters retain the truth value of
their class.

Two simultaneous inductions now do the work.  First, every original
\(\alpha\)-edge maps to a filtered \(\alpha\)-edge.  Second, whenever
\([\alpha]B\in\Delta\), truth of \([\alpha]B\) in a representative propagates
\(B\) across every filtered \(\alpha\)-edge.  Union and composition use their
displayed closure entries, and a test edge is an identity edge whose source
satisfies its test formula.  For the only delicate case, suppose
\([\beta^{*}]B\in\Delta\).  Closure also contains
\([\beta][\beta^{*}]B\).  Along a finite filtered path
\[
 [s_0]R_\beta^{f}[s_1]R_\beta^{f}\cdots R_\beta^{f}[s_n],
\]
the second induction propagates \([\beta^{*}]B\) from each class to the next;
the first class has it by hypothesis, while star unfolding gives both \(B\)
there and the one-step obligation needed to continue.  Repeating this argument
gives \(B\) at every class reached.  This proves the forward box direction.  For the reverse
direction, failure of \([\alpha]B\) in the original model supplies an actual
\(\alpha\)-successor where \(B\) fails, and the first induction maps that
witness to the quotient.  Structural induction on \(B\in\Delta\) therefore
gives the truth lemma
\[
 M,s\models B
 \quad\Longleftrightarrow\quad
 M^{f},[s]\models B.
\]
The class of the original satisfying state satisfies \(C\), so every
satisfiable formula has a model with at most
\(2^{|\operatorname{FL}(C)|}\) states.

Fischer and Ladner then search this finite space.  Under their coding and input
length \(n\), the paper states a nondeterministic upper bound \(c^n\) and a
deterministic lower bound of the form \(d^n/\log n\), for constants
\(c,d>1\).  Enumerating every exponential-size model deterministically would
instead give a double-exponential procedure.  The deterministic
single-exponential type-deletion algorithm and the exact
\(\mathsf{EXPTIME}\)-complete classification below are later results and are
kept under their later source \cite[pp.~199--211]{FischerLadner1979}.

\subsubsection{Kozen--Parikh's exact completeness boundary}

Kozen and Parikh use the test-free grammars
\[
 \alpha::=a\mid\alpha\cup\beta\mid\alpha;\beta\mid\alpha^{*},
 \qquad
 X::=p\mid X\lor Y\mid\neg X\mid\langle\alpha\rangle X,
\]
and define \([\alpha]X:=\neg\langle\alpha\rangle\neg X\).  Tests and converse
are explicitly omitted.  In the paper's diamond notation, the seven
Segerberg axiom groups are
\[
\begin{array}{rll}
 (1)&\text{all propositional tautologies},\\
 (2)&\langle\alpha\rangle\bot\leftrightarrow\bot,\\
 (3)&\langle\alpha\cup\beta\rangle X
       \leftrightarrow
       (\langle\alpha\rangle X\lor\langle\beta\rangle X),\\
 (4)&\langle\alpha\rangle(X\lor Y)
       \leftrightarrow
       (\langle\alpha\rangle X\lor\langle\alpha\rangle Y),\\
 (5)&\langle\alpha;\beta\rangle X
       \leftrightarrow\langle\alpha\rangle\langle\beta\rangle X,\\
 (6)&\langle\alpha^{*}\rangle X
       \leftrightarrow
       (X\lor\langle\alpha\rangle\langle\alpha^{*}\rangle X),\\
 (7)&\langle\alpha^{*}\rangle X
       \to
       X\lor\langle\alpha^{*}\rangle
          (\neg X\land\langle\alpha\rangle X).
\end{array}
\]
The rules are modus ponens and theorem-only generalization
\(X/[\alpha]X\) \cite[pp.~113--114]{KozenParikh1981}.

The two star axioms are sound for different reasons.  Since
\(R_{\alpha^{*}}\) is the identity relation together with one
\(R_{\alpha}\)-step followed by an \(R_{\alpha^{*}}\)-path, (6) is just the
relational equation
\[
 R_{\alpha^{*}}=I\cup(R_{\alpha}\mathbin{;}R_{\alpha^{*}}).
\]
For (7), suppose a finite \(\alpha\)-path from \(s=s_0\) reaches an
\(X\)-state.  If \(s\models X\), its first disjunct holds.  Otherwise choose the
least \(j>0\) with \(s_j\models X\).  Then \(s_{j-1}\models\neg X\),
\(s_{j-1}\models\langle\alpha\rangle X\), and \(s_{j-1}\) is
\(\alpha^{*}\)-reachable from \(s\), proving the second disjunct.  Dually this is
the familiar box induction principle
\[
 X\land[\alpha^{*}](X\to[\alpha]X)\to[\alpha^{*}]X.
\]

The completeness proof is finite and its three lemmas can be stated without
hiding the star step.  Begin with a consistent formula \(W\).  The paper's
test-free closure is the least set containing \(W\) and satisfying
\[
\begin{array}{ll}
 X\lor Y\in\operatorname{FL}(W)&\Rightarrow X,Y\in\operatorname{FL}(W),\\
 \neg X\in\operatorname{FL}(W)&\Rightarrow X\in\operatorname{FL}(W),\\
 \langle\alpha\rangle X\in\operatorname{FL}(W)&\Rightarrow X\in\operatorname{FL}(W),\\
 \langle\alpha\cup\beta\rangle X\in\operatorname{FL}(W)
   &\Rightarrow\langle\alpha\rangle X,\langle\beta\rangle X
     \in\operatorname{FL}(W),\\
 \langle\alpha;\beta\rangle X\in\operatorname{FL}(W)
   &\Rightarrow\langle\alpha\rangle\langle\beta\rangle X
     \in\operatorname{FL}(W),\\
 \langle\alpha^{*}\rangle X\in\operatorname{FL}(W)
   &\Rightarrow\langle\alpha\rangle X,
     \langle\alpha\rangle\langle\alpha^{*}\rangle X
     \in\operatorname{FL}(W).
\end{array}
\]
If this closure is \(\{X_1,\ldots,X_n\}\), an \emph{atom} is a consistent
conjunction \(Y_1\land\cdots\land Y_n\), with each
\(Y_i\in\{X_i,\neg X_i\}\).  This is the source's Boolean-algebra term, not a
new kind of program state.  Write \(C\leq X\) for
\(\derives C\to X\).  The canonical states are the atoms; for a primitive
program \(a\), put
\[
 C R_a D
 \quad\Longleftrightarrow\quad
 C\land\langle a\rangle D\text{ is consistent},
\]
and extend \(R\) by union, composition, and reflexive-transitive closure.

\paragraph{Lemma 1: consistent diamonds yield program edges.}
For every program \(\alpha\), consistency of
\(C\land\langle\alpha\rangle D\) implies \(C R_\alpha D\).  Choice follows
from (3).  For composition, write truth as the finite disjunction of all atoms;
axioms (4) and (5) produce an intermediate atom \(E\) for which both
\(C\land\langle\beta\rangle E\) and
\(E\land\langle\gamma\rangle D\) are consistent, and the induction hypothesis
gives \(C R_\beta E R_\gamma D\).

For \(\alpha=\beta^{*}\), let \(G\) be the least set of atoms containing \(C\)
and closed under every consistent one-step demand
\(E\land\langle\beta\rangle F\).  The program induction hypothesis makes every
member of \(G\) \(\beta^{*}\)-reachable from \(C\).  Let
\(Y=\bigvee_{E\in G}E\), a finite formula.  By closure of \(G\),
\(Y\land\langle\beta\rangle\neg Y\) is inconsistent, so
\(\derives Y\to[\beta]Y\).  Generalization and the induction axiom yield
\(\derives Y\to[\beta^{*}]Y\).  If \(D\notin G\), then
\(D\leq\neg Y\), making
\(C\land\langle\beta^{*}\rangle D\) inconsistent.  Hence \(D\in G\), and
\(C R_{\beta^{*}}D\).  This is the only place in Lemma~1 where the induction
axiom is used.

\paragraph{Lemma 2: modal witnesses are exactly canonical successors.}
For every \(\langle\alpha\rangle X\in\operatorname{FL}(W)\),
\[
 C\leq\langle\alpha\rangle X
 \quad\Longleftrightarrow\quad
 \exists D\,(C R_\alpha D\land D\leq X).
\]
For the forward direction, \(X\) is provably the finite disjunction of the atoms
below it; diamond distribution selects a consistent
\(C\land\langle\alpha\rangle D\), and Lemma~1 supplies the edge.  The reverse
direction is an induction on \(\alpha\).  Choice and composition use their
closure clauses.  In the atomic base, if \(C\not\leq\langle a\rangle X\), the
complete signed atom \(C\) entails its negation, contradicting consistency of
\(C\land\langle a\rangle D\) and \(D\leq X\).  In the star case, take a
finite \(\beta\)-path to \(D\), start
with \(D\leq X\to\langle\beta^{*}\rangle X\), and propagate
\(\langle\beta^{*}\rangle X\) backwards one edge at a time using the closure
formula \(\langle\beta\rangle\langle\beta^{*}\rangle X\) and axiom (6).

\paragraph{Lemma 3: truth lemma.}
For \(X\in\operatorname{FL}(W)\),
\[
 C\models X\quad\Longleftrightarrow\quad C\leq X.
\]
The propositional cases follow from the complete signed choice in each atom;
the diamond case is exactly Lemma~2 plus the formula induction hypothesis.
Finally, because \(W\) is consistent, some atom \(C\) satisfies \(C\leq W\),
and Lemma~3 gives \(C\models W\).  Thus every consistent test-free formula has
a finite model, and every valid test-free formula is derivable
\cite[pp.~114--118]{KozenParikh1981}.

Nothing in this proof licenses inserting a test \(B?\) into the 1981 grammar.
The next source adds the test axiom, the test closure case, and the corresponding
model argument explicitly.

\subsubsection{Harel--Kozen--Tiuryn's test-containing calculus and exact bounds}

Harel, Kozen, and Tiuryn's Axiom System 5.5 is the exact test-containing
controller for the registered sheet.  It adds
\([B?]A\leftrightarrow(B\to A)\) and prints the box versions of choice,
composition, star unfolding, and induction.  The sheet does not repeat the
source's conjunction-distribution axiom as primitive because it follows from
\(K\), necessitation, and propositional reasoning: necessitate
\(A\to(B\to A\land B)\) for one direction, and necessitate the two projections
for the other.  Its curried star-induction axiom is propositionally equivalent
to the source's conjunctive form.  The registered calculus therefore has the
strength of Axiom System 5.5 rather than a silently weakened fragment
\cite[\S5.5]{HarelKozenTiuryn2000}.

Tests can be seen directly in a three-state model.  Let
\[
 S=\{s,t,u\},\qquad
 R_a=\{(s,t),(s,u)\},\qquad
 V(p)=V(q)=\{t\}.
\]
Then \(R_{p?}=\{(t,t)\}\), so
\(R_{a;p?}=\{(s,t)\}\).  Consequently
\[
 s\models\langle a;p?\rangle q,
 \qquad
 s\models[a;p?]q,
 \qquad
 s\not\models[a]p.
\]
The diamond has the surviving \(t\)-witness; the box quantifies only over that
surviving witness; the untested box fails because \(u\) is an \(a\)-successor
that does not satisfy \(p\).  At a state with no
\(a;p?\)-successor, the corresponding box would be vacuously true while the
diamond would be false.

Section 6.1 supplies the simultaneous closure recursion used above.  Lemma 6.4
proves filtration, and Theorem 6.5 gives a model with at most \(2^{|A|}\) states
under the book's symbol-count convention.  Theorem 7.6 obtains deductive
completeness by filtering the universal nonstandard model; this is the
test-containing completion of the source route, not a restatement of the
test-free 1981 theorem \cite[\S\S6.1--6.2,7.1]{HarelKozenTiuryn2000}.

Algorithm 8.2 turns the finite vocabulary into a deterministic
single-exponential procedure.  For \(\Delta=\operatorname{FL}(A)\), it proceeds
as follows.
\begin{enumerate}
\item Form \(M_0\), the at most \(2^{|\Delta|}\) complete signed assignments to
  \(\Delta\): for each \(B\in\Delta\), exactly one of \(B,\neg B\) occurs.
\item Delete every assignment that fails a local propositional, choice,
  composition, test, or star-unfolding axiom.  The survivors form \(M_1\).
  Atomic edges satisfy
  \[
   uR_av
   \quad\Longleftrightarrow\quad
   \forall[a]B\in\Delta\,([a]B\in u\Longrightarrow B\in v),
  \]
  and compound relations are evaluated compositionally.
\item Repeatedly delete a state \(u\) containing
  \(\neg[\alpha]B\) when no current \(\alpha\)-successor \(v\) contains
  \(\neg B\).  The source chooses a violating \(\alpha\) of minimum length;
  continue until no deletion is possible.
\end{enumerate}
The third step enforces every existential demand
\(\neg[\alpha]B\equiv\langle\alpha\rangle\neg B\).  Actual consistent
filtered types are never deleted.  Conversely, at the fixed point the local
constructor equations and realized diamond demands prove a truth lemma for
the surviving structure.  Thus \(A\) is satisfiable exactly when a surviving
type contains it.  There are exponentially many candidates, at most that many
deletions, and even scanning all candidate edges costs only
\(2^{O(|A|)}\); Theorem 8.4 therefore gives deterministic exponential time
\cite[\S8.1]{HarelKozenTiuryn2000}.

The matching lower bound has a separate role.  Theorem 8.5 reduces acceptance
for a polynomial-space alternating one-tape Turing machine to PDL
satisfiability.  One atomic program \(\mathsf{Next}\) links configurations;
indexed proposition letters record the control state, head position, and tape
symbol.  Local formulas enforce legal configurations and moves.  At an
existential control state, \(\langle\mathsf{Next}\rangle\mathsf{Accept}\)
requires one accepting successor; at a universal state,
\([\mathsf{Next}]\mathsf{Accept}\) requires all successors to accept.  Star
forces the local constraints at every reachable configuration:
\[
 \mathsf{Start}
 \land[\mathsf{Next}^{*}]
   (\mathsf{Config}\land\mathsf{Move}\land\mathsf{Acceptance})
 \land\mathsf{Accept}.
\]
The construction is polynomial in the machine/input description and is
satisfiable exactly when the machine accepts.  Alternating polynomial space is
deterministic exponential time, so this supplies hardness; Algorithm 8.2
supplies membership.  These two arguments, not the finite-model theorem alone,
establish \(\mathsf{EXPTIME}\)-completeness
\cite[\S8.2]{HarelKozenTiuryn2000}.

The base theorem must not be transferred unchanged to a nearby logic merely
because its name still contains \enquote{PDL}.  The source-separated boundaries
are:
\begin{center}
\begin{tabularx}{\textwidth}{@{}l X X@{}}
\toprule
System & Exact change from base PDL & Why it is a separate result \\
\midrule
PDL
  & Arbitrary binary atomic-program relations; union, composition, star, and
    tests by arbitrary formulas.
  & This is the language/model class of Axiom System 5.5 and Theorems 6.5,
    7.6, 8.4, and 8.5. \\
DPDL
  & The syntax is unchanged, but atomic programs are interpreted as partial
    functions.
  & The model class validates
    \(\langle a\rangle A\to[a]A\), which base PDL does not; completeness needs
    that additional axiom and a modified construction. \\
Test-free and poor-test PDL
  & PDL(0) forbids tests; poor-test PDL permits only Boolean combinations of
    proposition letters as tests.
  & The languages have different expressive strength.  PDL(0) omits the test
    axiom, although the source's EXPTIME lower bound already uses no tests. \\
CPDL
  & Adds program converse \(\alpha^{-}\), interpreted by reversing the binary
    relation.
  & Converse strictly increases expressive power and requires converse axioms
    in its completeness proof. \\
LPDL and RPDL
  & Add predicates asserting termination of all computations or
    well-foundedness of a program relation.
  & Those properties are not expressible in base PDL, and the base filtration
    lemma fails for these extensions. \\
Concurrent PDL
  & Adds a parallel program constructor and interprets a program by a relation
    from one state to a set of possible final states.
  & The denotation is no longer an ordinary binary relation; its semantics and
    completeness axiom are separately proved. \\
\bottomrule
\end{tabularx}
\end{center}
These distinctions are the organization of Chapter 10, not terminological
variations on one inherited theorem \cite[ch.~10]{HarelKozenTiuryn2000}.
